\documentclass{article}
\usepackage{amsmath,amssymb,amsthm}
\newtheorem{lemma}{Lemma}\newtheorem{theorem}{Theorem}
\begin{document}
\title{Excluding the mixed zero-slack parity profile by common low-face geometry}
\maketitle
\section{M01: Exact original domain and imported interface}
Let $G\subseteq Q_8$ be a simple C4-free spanning graph with683 edges.
Suppose either one of its actual parity classes has degree counts
\[
(n_0,\ldots,n_8)=(0,0,1,2,7,61,57,0,0).                 \tag{1}
\]
There is no condition on the opposite parity. We use the actual-face
identity and nonnegative-loss theorem in
\texttt{../q8\_equality\_profiles/main.tex}, equivalently the stronger
\texttt{../q8\_crowded\_faces/main.tex}. Its historical local316/317/321
imports retain their recorded scope; they are not replayed here.
For(1), $C=333+2\cdot177+7\cdot66=1149$, so every one of the112
actual coordinate $Q_6$ faces has exactly zero or two specified-side
roots of restricted degree at most two. This is the sole imported
mathematical consequence needed below.

Denote the unique global degree-two root by$z$, the degree-three roots
by$x_A,x_B$ with actual retained direction masks$A,B$, and the seven
 degree-four roots by$y_j$ with masks$R_j$. Thus$|A|=|B|=3$ and
$|R_j|=4$. Every other specified root has degree at least five and
remains of degree at least three after two directions are omitted.
All ten named roots are distinct and have the same parity.

One common cube translation sends$z$ to0 and all ten roots to even
vertices. A common coordinate permutation, with an optional exchange
of$x_A,x_B$, makes
\[
 A=\{0,1,2\},\qquad
 B=\{0,\ldots,t-1\}\cup\{3,\ldots,5-t\},\quad t=0,1,2,3. \tag{2}
\]
The interval is empty when appropriate; at$t=3$, $B=A$.
This transports every position, retained mask and fixed-bit face
simultaneously. We do not normalize any other root independently.
We do not constrain the two retained directions at$z$: dropping them,
and all square-freeness restrictions, will already give a contradiction.

\section{M02: A complete finite direction domain, allowing repetitions}
For an omitted coordinate pair$I=\{a,b\}$, the low roots across its
four fixed-bit faces are exactly$z$, $x_A$ if$I\cap A\ne\varnothing$,
$x_B$ if$I\cap B\ne\varnothing$, and$y_j$ if$I\subseteq R_j$.
The four individual low counts are0 or2. Summing them modulo two gives
\[
 \bigoplus_{j=1}^7 k(R_j)
 =\mathbf1\oplus k([8]\setminus A)\oplus k([8]\setminus B),       \tag{3}
\]
where$k(T)$ is the28-bit incidence vector of the coordinate pairs
contained in$T$. Repeated$R_j$ are allowed in(3).

Any multiset of seven masks has an odd-multiplicity support of size
$r\in\{1,3,5,7\}$. Conversely, every possible seven-mask multiset
maps to one such support; all other copies occur in identical pairs
and cancel in(3). There are exactly70 possible four-masks.
The finite check in\texttt{classify\_directions.py} exhausts every odd
support: size1 and3 directly; size5 by the disjoint split1+4;
size7 by the disjoint split3+4. All quadruples and triples are generated
by combinations of the full70-mask list. Exact bitwise XOR is used,
and every matching union is retained as a set. Hence these joins have
no probabilistic or solver completeness assumption. The exact counts are
\[
\begin{array}{c|rrrr}
t& r=1&r=3&r=5&r=7\\\hline
0&0&0&0&0\\1&0&0&0&0\\2&0&0&0&28\\3&0&0&0&240
\end{array}                                                   \tag{4}
\]
Thus$t=0,1$ are impossible; for$t=2,3$ the seven masks are distinct.
This excludes repetitions by exhaustive parity support, rather than
assuming them absent.

For each of the latter two cases, enumerate all40320 coordinate
permutations and keep those preserving the unordered pair$\{A,B\}$.
Their orders are96 and720. Apply each such permutation to all seven
masks together and relabel the seven degree-four roots accordingly.
The resulting direction orbit representatives, as decimal bitmasks,
are
\[
\begin{array}{c|l|r}
t&\{R_1,\ldots,R_7\}&\text{orbit size}\\\hline
2&15,51,83,99,147,163,195&1\\
2&15,51,85,101,149,165,195&12\\
2&15,51,92,108,156,172,195&3\\
2&15,53,90,108,156,170,197&12\\
3&15,51,85,105,153,165,195&90\\
3&15,57,90,108,156,170,201&30\\
3&27,45,78,120,156,170,201&120
\end{array}                                                    \tag{5}
\]
\texttt{generate\_orbits.py} checks membership of every orbit in the
full saved domain and that the disjoint orbits have sizes28 and240.
\texttt{check\_directions\_independent.py} reconstructs the direction
joins in the reverse indexing order and independently checks that
these coordinate-permutation orbits cover the same complete domain.
No assertion is made that the abstract Gram rank alone gives(4).

\section{M03: Actual fixed-bit position exhaustion}
For each row of(5), the nine unknown actual positions range over
\[
                  E^*=\{v\in\{0,1\}^8:\mathrm{wt}(v)
                                      \text{ is even},\ v\ne0\}, \tag{6}
\]
which has127 elements. They must be pairwise distinct. For every
omitted pair$I$, let$T_I$ be the active roots listed in M02, including
$z=0$. Partition these roots by their actual two fixed bits$v_I$.
The exact remaining requirement is
\[
              |\{v\in T_I:v_I=c\}|\in\{0,2\}
                    \quad(c\in\{00,01,10,11\}).                 \tag{7}
\]
These are the112 actual-face conditions, not just28 direction totals.
The first row of(5) already has$|T_I|=10$ for some$I$, impossible
because four cells of capacity two hold at most eight roots.

Here is a finite exact algorithm for every other row.
For each$I$, enumerate all pairings of$T_I$. Give the pair containing
$z$ fixed-bit value00. Give the remaining pairs distinct values from
$01,10,11$, in every possible order. If$|T_I|>8$, reject immediately.
These relations are exactly(7): the occupied cells in a legal assignment
are its disjoint pairs, and conversely every such paired assignment
has every cell of size0 or2. There are no omitted relation tuples.

Give each root its full domain(6). Repeatedly, for each relation,
discard a tuple if some projected value is absent from the current
position domain of its root; remove from every root domain the positions
whose projection has no surviving tuple. Whenever a root domain is a
singleton, remove that position from all other root domains. Empty
domain or empty relation rejects a branch. These deletions are sound:
a completion using the current domains would use a relation tuple
whose every coordinate value is present, and could not reuse the
singleton position of a distinct root.

If propagation stabilizes and some root has several positions, choose
a root of least domain size and branch over every remaining position,
recursing with that root fixed. If every domain is a singleton, accept.
Induction on the number of unfixed positions proves this exhaustive
branching rejects a row only when no solution to(6)--(7) exists.
All domains are finite and every branch fixes a previously unfixed root.
The implemented checker has no cutoff, timeout or heuristic rejection.

The bitset implementation in
\texttt{check\_positions\_independent.py --allow-c4} returns EMPTY
for all seven rows, with search-node counts
\[
                       0,1,8,8,9,7233,11758.                    \tag{8}
\]
The maximum branching depths are respectively undefined,0,1,1,1,5,4.
The exact outputs are in\texttt{POSITION\_ALLOW\_C4.json}.
This run enforces distinctness but permits all C4s and omits even the
degree-two root's retained directions. Thus its domain contains every
actual input in M01. The earlier list-domain checker
\texttt{position\_gac.py} separately rejects all rows with square
conflicts present, using directly enumerated projected tuples;
its different final node counts are not used in(8).

\begin{lemma}
There do not exist ten distinct same-parity positions in$Q_8$, one
always active root, two roots active on omitted pairs meeting their
three-masks, and seven roots active on pairs contained in their
four-masks, such that all112 fixed-bit active counts are0 or2.
\end{lemma}
\begin{proof}
The one common normalization(2), complete odd-support check(4),
orbit cover(5), and exact actual-position exhaustion(6)--(8) cover
all inputs and reject every one. The arithmetic is finite binary and
integer arithmetic. The paired-relation construction and sound exhaustive
branching above are the certificate interpretation of the supplied code.
\end{proof}

\section{M04: Return to the original graph and precise assurance}
\begin{theorem}
No683-edge C4-free subgraph of$Q_8$ has the parity profile(1) on
either side, with the opposite side unrestricted.
\end{theorem}
\begin{proof}
An arbitrary such graph gives the ten distinct same-parity roots and
all112 zero-or-two actual low counts by M01. Their actual retained
masks and positions therefore satisfy the excluded finite geometric
lemma. This is a contradiction.
\end{proof}
The already-derived53 double-low faces,59 zero-low faces,12 common
$Q_3$ holes and ordinary/high-star completion obligations are not assumed
independently and need not be completed: a weaker necessary geometric
condition already has empty domain. This does not discard ordinary
roots by assuming an independent completion. No graph, partial marked
witness, or fractional relaxation is presented as a counterexample.

This is an author candidate computational proof with two locally written
implementations, not a fresh-agent proof review. Root should independently
review the complete-domain bridge and exact computations before accepting
the row exclusion. Historical imports retain their own assurance; no
historical SAT trace, Lean, or TeX/PDF compilation is claimed. Other
683-edge parity profiles remain open, and the unrestricted bound
$\operatorname{ex}(Q_8,C_4)\le682$ is not proved here.
\end{document}
